\section{Effect of the VLM and Human Audit}
\label{app:vlm-human-audit}

\paragraph{VLM Role and Leakage Control}
\textbf{Where VLM annotations enter.}
\model uses the VLM as a fixed attribute annotator. Given a task-specific set of attribute fields and closed-set label options, the VLM assigns a discrete value to each image. The selector then uses these values, class labels, and generated-data attribute counts to choose the source subset. The repair-slice indicator is the binary conjunction of the target attribute values for the corresponding task: red hair, brown skin, and sad emotion for CelebA; pink color and fabric texture for ImageNet10; brown clothes and rainy weather for COCO; and red object color and cloudy weather for KITTI. The main repair metrics in Tables~\ref{tab:main-fr} and~\ref{tab:app-main-imagenet}--\ref{tab:app-main-kitti} are standard task metrics computed from benchmark labels, so the VLM does not grade the final repaired model. Test-set attributes are used only for post-training slice diagnostics.

\vspace{0.5em}
\textbf{Leakage control.}
The attribute vocabulary, target fields, and VLM annotation prompt are fixed before running the source-selection methods. The assignment stage does not receive model predictions, failure rates, method names, selected-set membership, or downstream results. In settings where the same VLM family is used for target-slice wording and image-level assignment, the two roles are still separated: the slice description is fixed first, and assignment is then run as a frozen closed-set classifier over images. Source selection uses only source-training attributes and generated-pool attributes. Thus a possible failure mode is not test-label leakage, but inaccurate or self-consistent metadata. The audit below tests the target fields with independent human labels.

\paragraph{Human Audit Protocol}
We conduct a blinded human audit with five AI graduate students. Each annotator labels 100 CelebA images, 100 COCO images, and 100 KITTI images using the same closed-set options as the VLM. Annotators see only the image, dataset name, and sample identifier; they do not see VLM labels, target flags, model predictions, data split, or method outputs. We aggregate human labels by majority vote and compare them with the VLM labels at two levels: exact agreement on the audited fields and binary agreement on the repair-slice indicator after applying the target values for those fields. For example, the COCO slice indicator is true only when the audited image has brown clothes and rainy weather, and the KITTI indicator is true only when it has red object color and cloudy weather. For CelebA, the indicator uses the full repair conjunction over red hair, brown skin, and sad emotion, rather than a single field such as hair color alone. The binary agreement is the most relevant check for \model because the selector uses these attributes to identify repair-relevant source examples.

\begin{table}[h]
\centering
\caption{Human audit of VLM-assigned attributes. Five AI graduate students label each audited image using the same closed-set options as the VLM. Exact field agreement measures agreement on the audited attribute values, while slice-indicator agreement measures agreement on the binary target conjunction used for repair.}
\label{tab:vlm-human-audit}
\small
\setlength{\tabcolsep}{3pt}
\begin{tabular}{p{0.14\linewidth} C{0.27\linewidth} C{0.27\linewidth} C{0.22\linewidth}}
\toprule
\rowcolor{sectiongray} \textbf{Dataset} & \textbf{Target field 1 agree.} & \textbf{Target field 2 agree.} & \textbf{Slice-indicator agree.} \\
\midrule
CelebA & Hair color: 98\% & Skin tone: 99\% & 92\% \\
COCO & Clothes color: 97\% & Weather: 91\% & 90\% \\
KITTI & Object color: 98\% & Weather: 96\% & 94\% \\
\bottomrule
\end{tabular}
\end{table}

The audit is designed to measure agreement rather than to estimate the prevalence of rare target conjunctions. As shown in Table~\ref{tab:vlm-human-audit}, the VLM and human consensus usually agree on the audited target fields and on the resulting binary slice indicators. This supports the use of the VLM as a scalable attribute annotator rather than as an oracle: the VLM provides metadata, and the human audit checks that the target fields used for repair are visually grounded.

\paragraph{Sensitivity to attribute noise.}
We also test whether \model is brittle to noisy attributes. For each audited dataset and each target field, we independently corrupt a fraction $p$ of VLM labels by replacing the original value with a different valid value from the same field vocabulary. We then rerun source selection with the corrupted attributes while keeping the generated repair set, source ratio, and training recipe fixed. Table~\ref{tab:attribute-noise-sensitivity} reports the change from the clean-VLM \model result at source ratio $0.5$ for CelebA, COCO, and KITTI. Small changes under moderate noise would indicate that the selector depends on the aggregate attribute distribution rather than on a small number of individual VLM decisions.

\begin{table}[h]
\centering
\caption{Sensitivity to corrupted VLM attributes at source ratio $0.5$. Entries report the change from the clean-VLM \model result for each task metric.}
\label{tab:attribute-noise-sensitivity}
\small
\setlength{\tabcolsep}{5pt}
\begin{tabular}{l c c c c}
\toprule
\rowcolor{sectiongray} \textbf{Dataset} & \textbf{Metric} & \textbf{5\% noise} & \textbf{10\% noise} & \textbf{20\% noise} \\
\midrule
CelebA & Test ACC $\Delta$ & 0.007 & 0.01 & 0.09 \\
COCO & Test AP $\Delta$ & 0.002 & 0.003 & 0.02 \\
KITTI & Test mAP50--95 $\Delta$ & 0.003 & 0.004 & 0.06 \\
\bottomrule
\end{tabular}
\end{table}

The results show that \model is relatively stable to light and moderate attribute noise. At $5\%$ and $10\%$ corruption, the metric changes are small across all three datasets, which suggests that the selector is not driven by a few isolated attribute assignments. The larger shifts at $20\%$ noise, especially on CelebA and KITTI, are expected: once many target-field labels are replaced, the rarity and frequency counts no longer describe the diagnosed repair slice reliably. In this high-noise regime, \model no longer behaves like a clean semantic selector. Its ranking is partly driven by corrupted attribute values, while the soft class-balance term still preserves label coverage, so the selected subset moves toward a mixture of noisy attribute sampling and soft stratified sampling. This explains why the results begin to approach the behavior of random or stratified source selection when the attribute metadata is heavily corrupted. COCO is less sensitive in this range, suggesting that its selected subset is more strongly supported by aggregate attribute statistics or by other correlated visual factors in the source pool.

